\documentclass[12pt]{article}
\usepackage[utf8]{inputenc}
\usepackage[T1]{fontenc}
\usepackage{hyperref}
\title{Taabwa Language Resources for K-12 Curriculum Development}
\author{Compiled by Claude Sonnet 5 (claude-sonnet-5) for Andr\'as Kornai}
\date{}
\begin{document}
\maketitle
\section{Introduction}
Taabwa (Ichitaabwa, Kitabwa; ISO 639-3 \texttt{tap}) is a Bantu language of the DR Congo and Zambia. The figures found were about 250,000 speakers in the DRC (1982) and 32,000 in Zambia (2010 census). Little material was found.
\section{Primary Educational Resources}
None found.
\section{Linguistic and Reference Materials}
\begin{itemize}
\item Wikipedia article, which lists several unpublished Congolese theses on Taabwa and Shila grammar and phonology, and the 1907 van Acker Kitabwa--French dictionary (Tervuren): \url{https://en.wikipedia.org/wiki/Taabwa_language}
\end{itemize}
\section{Religious and Community Resources}
\begin{itemize}
\item Kitaabwa New Testament (The Word for the World International, CC BY-SA 4.0), from eBible: \url{https://www.ebible.org/epub/tap.epub}
\item Scripture Earth index: \url{https://www.scriptureearth.org/00i-Scripture_Index.php?iso=tap}
\item Joshua Project, Tabwa in Zambia: \url{https://www.joshuaproject.net/index.php/people_groups/15163/ZA}
\end{itemize}
\section{Digital and Audio Resources}
\begin{itemize}
\item Scripture Earth page listing Bible.is audio and video and Global Recordings Network audio: \url{https://www.scriptureearth.org/00eng.php?iso=tap}
\item African Proverbs, Sayings and Stories (Tangaza University), with a Taabwa e-book: \url{https://afriprov.tangaza.ac.ke/wp-content/uploads/2008/11/ebooks_etoka_taabwa.pdf}
\end{itemize}
\section{Curriculum Development Context}
Taabwa is not known to be taught in schools. The openly licensed New Testament and the proverb collection could serve as source texts. Primers and an orthography would need to be created with the community.
\end{document}